\section{Uniform moments and a positive-excess window near \(n^5\)}
\label{um-section}

Retain the actual roots \(a_k=3k\), \(w_k=a_k+\zeta\),
\(Q_k=a_k^2+a_k+1\), \(\alpha_k=w_k/\bar w_k\), for integers
\(B\leq k<2B\). The subset \(\mathcal S_B\) consists of roots
whose norm has a prime factor above \(12B\).
Theorem~\ref{pn-private} proves, for the actual roots,
\[
 |\mathcal S_B|\geq B/2-O(B\log\log B/\log B),
\]
and proves multiplicative independence of all their ratios.
The following uniform estimate removes the former fixed-moment
restriction while retaining this proved domain.

\begin{lemma}[Tuple rigidity uniform in its length]\label{um-rigidity}
Let \(B\geq5\) be an integer and \(q>8B\) a prime.
For every integer \(\nu\geq1\), two ordered \(\nu\)-tuples of
block roots, with unit conjugate denominators modulo \(q\),
whose ratio products agree modulo \(q^{2\nu}\), have exactly
equal algebraic ratio products in \(\mathbb Q(\zeta)\).
\end{lemma}
\begin{proof}
Write their actual element products as \(z=R+S\zeta\) and
\(z'=R'+S'\zeta\).
The norm coordinate bound and the finite telescoping expansion,
valid at every finite \(\nu\), give
\[
 |R|\leq2(7B)^\nu,\qquad |S|\leq2\nu(7B)^{\nu-1}.
\]
Thus for \(D=RS'-R'S\),
\[
 |D|\leq8\nu\,7^{2\nu-1}B^{2\nu-1}.
\]
The conjugate cross difference implies \(q^{2\nu}\mid D\).
Put \(x=49/64\). Since \(q>8B\), the height bound divided by
\(q^{2\nu}\) is less than
\[
 \frac8{7B}\nu x^\nu
 <\frac8{7B}\sum_{j=0}^{\nu-1}x^j
 <\frac8{7B}\frac1{1-x}
 =\frac{512}{105B}<\frac5B\leq1.
\]
Therefore \(D=0\), and \(z\bar z'=z'\bar z\) over the number
field. Their ratios agree. The products are nonzero; no sign
condition on \(R,S\) or bound on \(\nu\) is needed.
The hypothesis \(q>8B\) is used explicitly and is not replaced
by the earlier smaller cutoff.
\end{proof}

\begin{theorem}[A multiset bound at every precision]\label{um-multisets}
Let \(n>324\) be prime and \(B=n^4\). For a prime \(q>8B\)
and a positive integer \(\nu\), put
\[
 M_\nu(q)=\#\{k\in\mathcal S_B:q^{2\nu}\mid T_n(k)\}.
\]
If \(M_\nu(q)>0\), then
\begin{equation}\label{um-binomial}
 \binom{M_\nu(q)+\nu-1}{\nu}\leq3n.
\end{equation}
For an empty set the number of positive-length multisets is zero.
In particular, with \(r=\lceil\sqrt n\rceil\),
\[
 M_2(q)\leq\sqrt{6n},\qquad M_r(q)\leq3.
\]
\end{theorem}
\begin{proof}
Every hit gives a norm-one element killed by \(3n\) modulo
\(q^{2\nu}\). The split and inert simple-lifting proof of
Lemma~\ref{mc-torsion} bounds the group by \(3n\), since \(q>n\).
Every multiset of \(\nu\) roots gives one group product.
Equal images imply equality of actual ratio products by
Lemma~\ref{um-rigidity}. Independence from
Theorem~\ref{pn-private} then recovers all index multiplicities.
The multiset map is injective. Stars and bars gives
\eqref{um-binomial}, keeping repeated entries.

At \(\nu=2\), the inequality \(M(M+1)/2\leq3n\) proves the
first bound. If \(M_r\geq4\), the multisets on just four roots
would instead give
\[
 3n\geq\binom{r+3}{3}
      =\frac{(r+1)(r+2)(r+3)}6
      >\frac{r^3}6
      \geq\frac{n^{3/2}}6>3n.
\]
The last inequality uses \(\sqrt n>18\), a contradiction.
Here \(r\) grows with \(n\); the determinant estimate was first
proved uniformly for all finite lengths.
\end{proof}

\begin{theorem}[An explicit complete excess mean]\label{um-mean}
For every prime \(n>324\), \(B=n^4\), and real \(8B<Z\leq B^2\),
put
\[
 F_{8B,Z}(k)=\sum_{8B<q\leq Z}(v_q(T_n(k))-3)_+\log q.
\]
Then
\begin{equation}\label{um-eighty}
 \frac1B\sum_{k\in\mathcal S_B}\frac{F_{8B,Z}(k)}{t_k}
       \leq\frac{80Z}{Bn\log(Z/(3n))}.
\end{equation}
\end{theorem}
\begin{proof}
The complete positive excess satisfies
\((e-3)_+=\sum_{j\geq4}{\bf1}_{e\geq j}\).
Fix \(q>8B\), and put \(r=\lceil\sqrt n\rceil\).
Every layer \(4\leq j<2r\) has at most \(\sqrt{6n}\) roots
in \(\mathcal S_B\). There are at most \(2r\) layers, with
\(r\leq2\sqrt n\), so their contribution is at most
\[
 4\sqrt6\,n\log q<10n\log q.
\]
Every higher layer has at most three roots.
Writing \(L=\log(6B+1)\), the actual height bound
\(\log T_n(k)\leq3nL\) caps its depth by \(3nL/\log q\).
All higher layers together contribute at most \(9nL\).
Consequently
\begin{equation}\label{um-prime-ledger}
 \sum_{k\in\mathcal S_B}(v_q(T_n(k))-3)_+\log q
       \leq10n\log q+9nL.
\end{equation}
This includes empty layer ranges and pays every positive depth.

For \(q\leq Z\leq B^2\) and \(B\geq7\),
\(\log q\leq2\log B\) and \(L\leq2\log B\).
Also \(t_k\geq n\log(3B)\geq n\log B\). Thus the normalized
contribution at one prime is at most \(38/B\).
Positive excess excludes rank one: \(q>n\) gives no exponent
lifting, and \(q^2\) cannot divide either coprime old factor
\(a_k,a_k+1<6B+1\). Hence the exact rank is \(n\), and
\(q\equiv\pm1\pmod{3n}\).
The two-progression Brun--Titchmarsh estimate bounds the number
of eligible primes up to \(Z\) by
\[
 \frac{2Z}{(n-1)\log(Z/(3n))}.
\]
Multiplying by \(38/B\) and using \(76n/(n-1)<80\) proves
\eqref{um-eighty}. The normalization is by the whole block \(B\).
\end{proof}

\begin{corollary}[A logarithmic expansion beyond \(n^5\)]
\label{um-domain}
Fix \(\delta>0\). For all sufficiently large prime \(n\), put
\[
 Z_n=\left\lfloor Bn(\log n)^{1-\delta}\right\rfloor,\qquad B=n^4.
\]
There is an actual subset \(\mathcal P_{n,\delta}\) of relative
size at least
\[
 \frac12-
 O_\delta\!\left(\frac{\log\log n}{\log n}
                  +(\log n)^{-\delta/2}+(\log n)^{-1/2}\right)
\]
on which
\[
 \frac{F_{8B,Z_n}(k)}{t_k}\leq(\log n)^{-\delta/2},\qquad
 \frac{L_{8B}(T_n(k))}{t_k}\leq(\log n)^{-1/2}.
\]
\end{corollary}
\begin{proof}
For fixed \(\delta\), \(8B<Z_n\leq B^2\) eventually, and
\[
 \log(Z_n/(3n))
   =4\log n+(1-\delta)\log\log n-\log3+o(1)
   \sim4\log n.
\]
Theorem~\ref{um-mean} gives mean \(O_\delta((\log n)^{-\delta})\).
Markov at \((\log n)^{-\delta/2}\) gives its displayed error.
The full-mass estimate in Theorem~\ref{fml-mean}, with
Section~\ref{ebi-section}'s elementary inputs, at cutoff \(8B\)
has mean \(O(1/\log n)\).
Markov at \((\log n)^{-1/2}\) gives the other error. Intersect
these sets with Theorem~\ref{pn-private}.
In particular the interval \((B,8B]\) is paid within the full
mass up to \(8B\), without imposing a new cap on its depths.
\end{proof}

Let \(W_{Z_n}(k)\) denote the signed contribution above \(Z_n\).
The exact prime decomposition gives on this set
\begin{align}
 \frac{J(T_n(k))}{t_k}
 &\leq\frac{W_{Z_n}(k)}{t_k}
                   +(\log n)^{-\delta/2}+(\log n)^{-1/2},
       \label{um-signed}\\
 \frac{\log\operatorname{rad}(T_n(k))}{t_k}
 &\geq1-\frac4{3n}
       -\frac{(\log n)^{-\delta/2}+(\log n)^{-1/2}}3
       -\frac{(W_{Z_n}(k))_+}{3t_k}. \label{um-radical}
\end{align}
The second inequality uses the actual elementary angular estimate.
Taking \(\delta=1\) reaches \(Z_n=n^5\); any fixed \(0<\delta<1\)
gives the stated logarithmic expansion.
The complete positive excess in this interval is controlled on
the proved domain. Its full mass need not be small, and no size
of compensating negative credit is presumed.

The fifteen declarations of
\texttt{UniformMoment\allowbreak Arithmetic.lean} formalize the
numeric binomial gate, an explicit integer square-root threshold,
the full finite-depth list ledger and its real weighted versions,
and the constants \(38\) and \(80\) under explicit interfaces.
The list count equals the actual filtered-list length.
The threshold \(\lfloor\sqrt n\rfloor+1\) used in the code has
the same required bounds and also permits square \(n\) in the
numeric statements. Finite-ring construction, ideal valuations,
the uniform tuple height estimate and the analytic estimates are
not included in those fifteen formal claims.
The unbounded farther signed cost, the complementary roots whose
norms are smooth at this threshold, and the remaining exceptional
subsets are still unresolved.
